\documentclass[10pt]{beamer}

\usetheme{metropolis}
\definecolor{WiLabRed}{RGB}{197,18,48}
\setbeamercolor{frametitle}{fg=white,bg=WiLabRed}
\setbeamercolor{progress bar}{fg=WiLabRed!90}
\setbeamercolor{title separator}{fg=WiLabRed!90}
\setbeamercolor{progress bar in section page}{fg=WiLabRed!90}
\setbeamercolor{background canvas}{bg=white}
\setbeamercolor{alerted text}{fg=WiLabRed!90}

\usepackage{appendixnumberbeamer}

\usepackage{booktabs}
\usepackage[scale=2]{ccicons}

\usepackage{pgfplots}
\usepgfplotslibrary{dateplot}

\usepackage{xspace}
\newcommand{\themename}{\textbf{\textsc{metropolis}}\xspace}

\usepackage{marvosym}
%\usepackage{subfig}
\usepackage{graphicx}\graphicspath{{images/}}
\usepackage{subcaption}
\usepackage[framed]{./support/mcode}
\usepackage{listings}

\title{Lecture 09: Phase Shift Keying}
\subtitle{\textit{Software Defined Radio for Engineers} (Collins~\textit{et al.}), \textsection{4.2.4}}
\date{}
\author{\textbf{Alexander M. Wyglinski, Ph.D.}}
\institute{ \vspace*{1in}\hfill\includegraphics[height=1.125cm]{wilab_logo-A70916.eps} \qquad \includegraphics[height=1.125cm]{WPI_Inst_Prim_FulClr.eps}}
% \titlegraphic{\hfill\includegraphics[height=1.5cm]{logo.pdf}}

% Foot for all slides
\setbeamertemplate{frame footer}{\tiny \copyright~2018 by Alexander Wyglinski. This work is licensed under the Creative Commons Attribution-ShareAlike 4.0 International License. To view a copy of this license, visit http://creativecommons.org/licenses/by-sa/4.0/.}

\begin{document}

%\captionsetup[subfigure]{labelformat=empty}

%%%%%%%%%%%%%%%%%%%%%%%%%%%%%%%%%%%%%%%%%%%%%%%%%%%%%%%%%%

\maketitle




%%%%%%%%%%%%%%%%%%%%%%%%%%%%%%%%%%%%%%%%%%%%%%%%%%%%%%%%%%

\frame
{
  \frametitle{BPSK Modulation}

    \begin{itemize}
        \item Modulation Rule:
        \begin{equation}
        \begin{split}
            ``1"\rightarrow{s_{1}(t)}&=A\cdot\cos(\omega_{c}t+\theta)\\
            ``0"\rightarrow{s_{2}(t)}&=-A\cdot\cos(\omega_{c}t+\theta)\\
            &=A\cdot\cos(\omega_{c}(t)+\theta+\pi)\\
            &=-s_{1}(t)
        \end{split}
        \end{equation}
        \item As we have seen in the past few slides, determining $\varepsilon_{P}$ is important
        \begin{itemize}
        \item We will now show that $\varepsilon_{p}=4$ is the {\bf optimum}
        \end{itemize}
    \end{itemize}


}

\frame
{
  \frametitle{Solving for $\varepsilon_{p,\mathsf{BPSK}}$: Compute $d_{\min}^2$}

   \begin{itemize}
    \item Employ the definition for $d_{min}^2$ and solve
        \begin{equation}
        \begin{split}
            d_{\min}^2&=\int\limits_{0}^{T}(s_{1}(t)-s_{2}(t))^2{dt}\\
            &=4A^2\int\limits_{0}^{T}\cos^2(\omega_ct+\theta)dt\\
            &=\frac{4A^2T}{2}+\frac{4A^2}{2}\int\limits_{0}^{T}\cos(2\omega_ct+2\theta)dt\\
            &=2A^2T\rightarrow\mathsf{double~frequency~term~eliminated}\nonumber
        \end{split}
        \end{equation}
    \end{itemize}

}

\frame
{
  \frametitle{Solving for $\varepsilon_{p,\mathsf{BPSK}}$: Compute $\bar{E}_{b}$}

   \begin{itemize}
   \item Employ the definition for $\bar{E}_{b}$ and solve
    \begin{equation}
    \begin{split}
        E_{s_1}&=\int\limits_0^Ts_1^2(t)dt=A^2\int\limits_0^T\cos^2(\omega_ct+\theta)dt\\
        &=\frac{A^2T}{2}+\frac{A^2}{2}\int\limits_{0}^{T}\cos(2\omega_ct+2\theta)dt\\
        &=\frac{A^2T}{2}\rightarrow\mathsf{double~frequency~term~eliminated}\\
        E_{s_2}&=\frac{A^2T}{2}\quad\therefore\quad\bar{E}_b=P(0)\cdot{E_{s_2}}+P(1)\cdot{E_{s_1}}=\frac{A^2T}{2}\nonumber
    \end{split}
    \end{equation}
    \end{itemize}
}


\frame
{
  \frametitle{Solving for $\varepsilon_{p,\mathsf{BPSK}}$: The Finale}

    \begin{itemize}
        \item Applying the definition for the power efficiency yields:
        \begin{equation}
            \varepsilon_{p,\mathsf{BPSK}} = \frac{d_{\min}^2}{\bar{E}_{b}}=4
        \end{equation}
        \item This is suppose to be the largest possible value for $\varepsilon_{p}$ for a modulation scheme employing all possible signal representations, i.e., $M=2^b$ waveforms
    \end{itemize}
}

\frame
{
  \frametitle{Alternative Approach to Computing $d_{\min}^2$}

  \begin{itemize}
        \item Another way to compute $d_{\min}^2$ is to use the concept of {\it correlation} such that:
        \begin{equation}
            d_{\min}^2=\int\limits_{0}^{T}(s_2(t)-s_1(t))^2dt=E_{s_1}+E_{s_2}-2\rho_{12}
        \end{equation}
        where:
        \begin{equation}
            E_{s_i}=\int\limits_{0}^{T}s_i^2(t)dt\quad\mathsf{and}\quad\rho_{12}=\int\limits_{0}^{T}s_1(t)s_2(t)dt\nonumber
        \end{equation}
        \item In order to get a {\bf large} $\varepsilon_{p}$, we need to maximize $d_{\min}^2$, which means we want $\rho_{12}<0$
        \begin{itemize}
            \item $E_{s_1}=E_{s_2}=E=A^2T/2\rightarrow{d_{\min}^2}=2(E-\rho_{12})\rightarrow{\rho_{12}=-E}$
        \end{itemize}
  \end{itemize}

}


\frame
{
  \frametitle{An Example}

    \begin{itemize}
        \item Show that for the following signal waveforms:
        \begin{equation}
        \begin{split}
            s_1(t)&=A\cdot\cos(\omega_ct+\theta)\\
            s_2(t)&=0\nonumber
        \end{split}
        \end{equation}
        the power efficiency is equal to $\varepsilon_{p}=2$
        \begin{itemize}
            \item This represents a 3~dB loss relative to BPSK!!
            \item $\mathsf{Loss}=10\cdot\log_{10}\left(\frac{\varepsilon_{p,\mathsf{BPSK}}}{\varepsilon_{p,\mathsf{other}}}\right)$
        \end{itemize}
        \item Show that for the following signal waveforms:
        \begin{equation}
        \begin{split}
            s_1(t)&=A\cdot\cos(\omega_ct+\theta)\\
            s_2(t)&=A\cdot\sin(\omega_ct+\theta)\nonumber
        \end{split}
        \end{equation}
        the power efficiency is equal to $\varepsilon_{p}=2$
    \end{itemize}


}

\frame
{
  \frametitle{QPSK Modulation Format}

    \begin{itemize}
        \item So far we have seen modulation schemes that consist of just one of two waveforms
        \begin{itemize}
            \item We will now expand our signal constellation repertoire to four distinct waveforms per modulation scheme
        \end{itemize}
        \item In QPSK modulation, a signal waveform possesses the following representation:
        \begin{equation}
        \begin{split}
            {s_{i}(t)}=&\pm{A}\cdot\cos(\omega_{c}t+\theta)\\
            &\pm{A}\cdot\sin(\omega_{c}t+\theta)
        \end{split}
        \end{equation}
    \end{itemize}

}

\frame
{
  \frametitle{Computing $\varepsilon_{p,\mathsf{QPSK}}$}

    \begin{itemize}
       \item Solving for $d_{\min}^2$, we obtain:
        \begin{equation}
            d_{\min}^2=\int\limits_{0}^{T}\Delta{s^2(t)}dt=2A^2T
        \end{equation}
       \item To find $\bar{E}_b$, we need to average over all the signals, which is equal to (assuming $E_{s_1}=E_{s_2}=E_{s_3}=E_{s_4}=A^2T$):
         \begin{equation}
            \bar{E}_b=\frac{(E_{s_1}+E_{s_2}+E_{s_3}+E_{s_4})/4}{\log_2(M)}=\frac{A^2T}{2}
         \end{equation}
       \item Solving for the power efficiency, we get:
       \begin{equation}
        \varepsilon_{p,\mathsf{QPSK}}=\frac{d_{\min}^2}{\bar{E}_b}=4
       \end{equation}
       which is the same as BPSK {\bf but} with 2 bits per symbol!!
    \end{itemize}

}

\frame
{
  \frametitle{MPSK Modulation Format}
  
  \begin{itemize}
    \item Constant distance of signal constellation point to origin
    \item Consists of $M$ equally spaced points on a circle
    \item General expression for an $M$-ary PSK waveform:
    \begin{equation}
        s_i(t)=A\cdot\cos\left(\omega_ct+\frac{2\pi{i}}{M}\right),~\mathsf{for}~i=0,1,2,\ldots,M-1
    \end{equation}
    \item There are advantages and disadvantages with this modulation scheme
    \begin{itemize}
        \item As $M$ increases, the spacing between signal constellation points decreases $\rightarrow$ error robustness decreases
        \item Having information encoded in the phase results in constant envelope modulation, which is:
        \begin{itemize}
            \item Good for non-linear power amplifiers
            \item Robust to amplitude distortion channels
        \end{itemize}
    \end{itemize}
  \end{itemize}
  
}
\frame
{
  \frametitle{Computing $\varepsilon_{p,\mathsf{M-PSK}}$}
  
  \begin{itemize}
    \item Given $s_1(t)=A\cdot\cos(\omega_ct)$ and $s_2(t)=A\cdot\cos(\omega_ct+2\pi/M)$
    \begin{itemize}
        \item Calculate $d_{\min}^2=E_{s_1}+E_{s_2}-2\rho_{12}$ where:
        \begin{equation}
            E_{s_i}=\int\limits_0^Ts_i^2(t)dt=\frac{A^2T}{2},~\mathsf{for}~i=1,2\quad
        \end{equation}
        and:
        \begin{equation}
            \rho_{12}=\int\limits_0^Ts_1(t)s_2(t)dt=\frac{A^2T}{2}\cos\left(\frac{2\pi}{M}\right)
        \end{equation}
        which yields $d_{\min}^2=A^2T(1-\cos(\frac{2\pi}{M}))$
    \end{itemize}
  \end{itemize}
  
  
}
\frame
{
  \frametitle{Computing $\varepsilon_{p,\mathsf{M-PSK}}$}
  
  \begin{itemize}
      \item Given $s_1(t)=A\cdot\cos(\omega_ct)$ and $s_2(t)=A\cdot\cos(\omega_ct+2\pi/M)$
    \begin{itemize} 
           \item The average bit energy $\bar{E}_b$ is equal to $\bar{E}_b=\frac{\bar{E}_s}{\log_{2}(M)}=\frac{\bar{E}_s}{b}$, where $\bar{E}_s=A^2T/2$
        \item Using the definition for the power efficiency, we see that $\varepsilon_{p,\mathsf{M-PSK}}=2b(1-\cos(\frac{2\pi}{M}))=4b\sin^2(\frac{\pi}{2^b})$
     \end{itemize}
  \end{itemize}
  
  
}

\frame
{
  \frametitle{Overall Comparison of $\varepsilon_{p}$}

    \begin{itemize}
        \item To determine how much power efficiency we are losing relative to $\varepsilon_{p,\mathsf{QPSK}}$, we use $\delta\mathsf{SNR}=10\cdot\log_{10}(\frac{\varepsilon_{p,\mathsf{QPSK}}}{\varepsilon_{p,\mathsf{other}}})$
    \end{itemize}
    \begin{table}[h]
    \caption{$\delta$SNR Values of Various Modulation Schemes.}
    \begin{tabular}{ccccc}
      \hline
      $M$ & $b$ & $M$-ASK & $M$-PSK & $M$-QAM \\
      \hline
      2  & 1 &  0    &  0    & 0 \\
      4  & 2 &  4    &  0    & 0 \\
      8  & 3 &  8.45 &  3.5  & (??) \\
      16 & 4 & 13.27 &  8.17 & 4.0 \\
      32 & 5 & 18.34 & 13.41 & (??) \\
      64 & 6 & 24.4  & 18.4  & 8.45 \\
      \hline
    \end{tabular}
    \end{table}

}
\frame
{
  \frametitle{Observations of $\varepsilon_{p}$}

    \begin{itemize}
        \item Two dimensional modulation is better than one dimensional modulation
        \item All modulation schemes studied are linear modulation schemes$\rightarrow$similar receiver complexity
    \end{itemize}
    
}

\end{document}
